\chapter{Dwarfism}

\section*{Definition and Overview}
Dwarfism is a term used to describe short stature in adults, usually defined as a height of less than 147 centimetres. There are more than 200 medical conditions that cause dwarfism, which are divided into two broad groups: proportionate short stature, in which the limbs and trunk are in proportion but small, and disproportionate short stature, in which the limbs are short relative to the trunk, as in achondroplasia. Most forms of dwarfism are genetic and present from birth or early childhood. This chapter explains the causes, health implications, and support available for people with dwarfism and their families.

\section*{Epidemiology}
Achondroplasia is the most common form of disproportionate dwarfism and affects approximately 1 in 25,000 to 30,000 births worldwide. Other forms of skeletal dysplasia are rarer. Proportionate short stature is more common and has many causes, including growth hormone deficiency and constitutional delay. Most people with dwarfism have normal intelligence and lead full, active lives, with the main challenges being medical complications, accessibility, and social attitudes.

\section*{Aetiology and Pathophysiology}
The causes of dwarfism are varied and are best understood through the two main categories.
\begin{itemize}
    \item \textbf{Disproportionate short stature (skeletal dysplasias):} conditions such as achondroplasia, in which bone and cartilage development is abnormal; achondroplasia is usually caused by a new mutation in the FGFR3 gene
    \item \textbf{Proportionate short stature:} growth hormone deficiency, hypothyroidism, chronic illness, malnutrition, or genetic syndromes affecting overall growth
    \item \textbf{Growth plate function:} in skeletal dysplasias, the growth plates in the long bones fail to produce bone normally, shortening the limbs
    \item \textbf{Inheritance:} some forms are inherited in dominant, recessive, or X-linked patterns, while many cases arise from new mutations in families with no history
    \item \textbf{Associated features:} depending on the cause, there may be effects on the spine, skull, joints, and internal organs
\end{itemize}

\section*{Clinical Features}
The features of dwarfism vary with the underlying condition.
\begin{itemize}
    \item Short stature, which may be evident at birth or become apparent in childhood when growth falls behind
    \item Short arms and legs relative to the trunk, with a large head and prominent forehead in achondroplasia
    \item Bowed legs, a curved lower spine, and a waddling gait
    \item Joint pain, stiffness, and reduced joint mobility
    \item Recurrent ear infections and hearing problems in achondroplasia
    \item Sleep apnoea, caused by a narrowed airway
    \item Compression of the spinal cord at the base of the skull or in the spine, causing neurological symptoms
    \item Delayed growth without other features in proportionate forms
\end{itemize}

\section*{Differential Diagnosis}
Determining the specific cause of short stature is important for management.
\begin{itemize}
    \item Constitutional delay of growth and puberty, a normal variant
    \item Familial short stature, in which parents are short
    \item Growth hormone deficiency and other endocrine disorders
    \item Hypothyroidism
    \item Chronic illness, including coeliac disease, inflammatory bowel disease, and kidney disease
    \item Genetic syndromes associated with short stature
    \item Skeletal dysplasias such as achondroplasia and other bone disorders
\end{itemize}

\section*{When to Seek Medical Care}
Children who fall well below the growth charts, whose growth rate slows, or who have features suggestive of a skeletal dysplasia should be assessed by a GP and referred to a paediatrician or clinical geneticist. Early assessment identifies treatable causes and allows families to access support.

\textbf{Call 000 (or local emergency services) for:}
\begin{itemize}
    \item Sudden weakness, numbness, or difficulty walking, which may indicate spinal cord compression
    \item Severe breathing difficulty or pauses in breathing
    \item New or worsening neurological symptoms such as weakness in the limbs or loss of bladder control
\end{itemize}

\section*{Diagnosis and Investigations}
Diagnosis aims to identify the specific cause of short stature.
\begin{itemize}
    \item \textbf{History and examination:} growth measurements over time, family history, and physical features
    \item \textbf{Growth charts:} plotting height, weight, and growth velocity against population standards
    \item \textbf{Blood tests:} growth hormone tests, thyroid function, and screening for chronic illness
    \item \textbf{Bone age X-ray:} X-ray of the hand and wrist to assess skeletal maturity
    \item \textbf{Genetic testing:} identifies mutations such as FGFR3 in achondroplasia
    \item \textbf{Imaging of the spine and skull:} to assess for compression in skeletal dysplasias
    \item \textbf{Specialist referral:} to endocrinologists, geneticists, and orthopaedic surgeons as needed
\end{itemize}

\section*{Management}
Management is tailored to the underlying cause and the individual's needs.
\begin{itemize}
    \item \textbf{Treatable causes:} growth hormone therapy for growth hormone deficiency, and thyroid replacement for hypothyroidism
    \item \textbf{Monitoring for complications:} regular review for ear infections, sleep apnoea, spinal compression, and joint problems
    \item \textbf{Orthopaedic care:} surgery for bowed legs, spinal problems, or nerve compression when needed
    \item \textbf{Sleep apnoea management:} assessment and treatment, including breathing support where required
    \item \textbf{Physiotherapy and occupational therapy:} to maintain mobility and independence
    \item \textbf{Educational support:} adjustments for any learning or hearing needs
    \item \textbf{Psychological support:} addressing the social and emotional aspects of short stature
    \item \textbf{Access and advocacy:} support for home and workplace modifications and community inclusion
\end{itemize}

\section*{Complications}
\begin{itemize}
    \item Spinal cord compression, which can cause weakness, numbness, and paralysis if untreated
    \item Sleep apnoea and respiratory problems
    \item Recurrent ear infections and hearing loss
    \item Joint pain and early arthritis
    \item Hydrocephalus, or fluid on the brain, in some forms
    \item Social isolation, bullying, and psychological distress
    \item Medical complications in pregnancy for women with skeletal dysplasias
\end{itemize}

\section*{Prognosis and Outcomes}
The outlook depends on the underlying cause. Most people with dwarfism have normal intelligence and a normal lifespan, and with appropriate medical care they can expect to live full, active lives. Complications such as spinal compression and sleep apnoea require ongoing surveillance, but are treatable. Supportive families, accessible environments, and inclusive communities make an important difference to quality of life.

\section*{Prevention and Self-Care}
\begin{itemize}
    \item Attend routine growth monitoring in childhood so problems are detected early
    \item Maintain regular follow-up with the specialist team managing the condition
    \item Report new neurological symptoms, breathing problems, or pain promptly
    \item Encourage activity within safe limits to maintain fitness and joint health
    \item Connect with support organisations for information, community, and advocacy
    \item Advocate for accessible environments in school, work, and public spaces
\end{itemize}

\section*{Sources and Further Reading}
The following reputable sources provide further information for healthcare students and patients seeking reliable, current advice about dwarfism.
\begin{itemize}
    \item Short Statured People of Australia. \textit{Support and advocacy for people with dwarfism.} https://www.sspa.org.au/
    \item NHS. \textit{Achondroplasia: overview and management.} https://www.nhs.uk/conditions/achondroplasia/
    \item Mayo Clinic. \textit{Dwarfism: symptoms and causes.} https://www.mayoclinic.org/diseases-conditions/dwarfism/
    \item MedlinePlus. \textit{Dwarfism.} https://medlineplus.gov/dwarfism.html
    \item Little People of America. \textit{Resources on dwarfism and skeletal dysplasia.} https://www.lpaonline.org/
    \item MSD Manual. \textit{Skeletal dysplasias: professional overview.} https://www.msdmanuals.com/
\end{itemize}
